\documentclass[12pt]{article}
\usepackage[a4paper, margin=1in]{geometry}
\usepackage{amsmath, amssymb, booktabs, setspace}
\setstretch{1.2}
\parskip=0.7em

\begin{document}


%============================================================
\section*{5. Factor Price Equalization and Real Incomes}



%-------------------------------------------------------------
\noindent\textbf{Solution and Explanation:}

\subsection*{(1) Stolper–Samuelson Theorem Recap}

\textbf{Definition:}  
The \textbf{Stolper–Samuelson theorem} states that an increase in the relative price of a good raises the real return to the factor used intensively in that good and lowers the real return to the other factor.

\noindent
Formally, if Computers (C) are capital-intensive:
\[
\frac{P_C}{P_T} \uparrow \; \Rightarrow \; r \uparrow, \; w \downarrow
\]
Conversely, if Textiles are labor-intensive:
\[
\frac{P_T}{P_C} \uparrow \; \Rightarrow \; w \uparrow, \; r \downarrow
\]

%-------------------------------------------------------------
\subsection*{(2) Quantitative Derivation}

We now derive how the percentage change in good prices affects factor prices.

\textbf{Step 1: Unit cost functions}

For each good produced with constant returns to scale and competitive markets, the price equals the unit cost:
\[
P_i = a_{Ki} r + a_{Li} w, \quad i = C, T
\]
where \(a_{Ki}\) and \(a_{Li}\) are input requirements per unit of output.

Taking percentage changes (denoted by hats):
\[
\hat{P}_i = \theta_{Ki}\hat{r} + \theta_{Li}\hat{w}
\]
where \(\theta_{Ki}\) and \(\theta_{Li}\) are cost shares of capital and labor in good \(i\), satisfying \(\theta_{Ki} + \theta_{Li} = 1\).

\textbf{Step 2: Set up system of equations}

For Computers (\(i=C\)) and Textiles (\(i=T\)):
\[
\begin{cases}
\hat{P}_C = \theta_{KC}\hat{r} + \theta_{LC}\hat{w} \\
\hat{P}_T = \theta_{KT}\hat{r} + \theta_{LT}\hat{w}
\end{cases}
\]

Using the given production exponents (which correspond to cost shares):
\[
\theta_{KC} = 0.6, \quad \theta_{LC} = 0.4; \qquad
\theta_{KT} = 0.3, \quad \theta_{LT} = 0.7
\]

\textbf{Step 3: Substitute known price changes}

In Canada:
\[
\hat{P}_C = +0.20, \quad \hat{P}_T = -0.10
\]

Substitute:
\[
\begin{cases}
0.20 = 0.6\hat{r} + 0.4\hat{w} \\
-0.10 = 0.3\hat{r} + 0.7\hat{w}
\end{cases}
\]

\textbf{Step 4: Solve the system for \(\hat{r}\) and \(\hat{w}\)}

Multiply the second equation by 2 to align the coefficients of \(\hat{r}\):
\[
\begin{cases}
0.20 = 0.6\hat{r} + 0.4\hat{w} \\
-0.20 = 0.6\hat{r} + 1.4\hat{w}
\end{cases}
\]

Subtract the second equation from the first:
\[
0.40 = -1.0\hat{w} \quad \Rightarrow \quad \hat{w} = -0.40
\]

Substitute into the first equation:
\[
0.20 = 0.6\hat{r} + 0.4(-0.40)
\Rightarrow 0.20 = 0.6\hat{r} - 0.16
\Rightarrow 0.36 = 0.6\hat{r}
\Rightarrow \hat{r} = 0.60
\]

\textbf{Interpretation:}
\[
\boxed{\hat{r} = +60\%, \quad \hat{w} = -40\%}
\]

\textbf{Conclusion:}  
A 20\% increase in the price of Computers (and a 10\% fall in Textiles) leads to:
- A 60\% increase in the real return to capital (\(r\))
- A 40\% fall in the real wage (\(w\))

This demonstrates that factor price changes are \emph{magnified} relative to goods price changes — a key result of the Stolper–Samuelson theorem.

\textbf{Economic meaning:}
- Capital owners in Canada gain substantially.
- Workers lose both relatively and absolutely (their real wage falls even though overall welfare may rise).


\end{document}
